% status: SKELETON
% Chapter 6 of the second edition. Plan: docs/STRUCTURE.md. Rules: AUTHORING.md.
\chapter{Napkin Math: Latency, Throughput and Cost per Token}\label{ch:napkin-math}

\begin{ChapterIntent}
The previous chapters give three numbers for any model on any machine:
bytes per token, FLOPs per token and cache bytes per token. This chapter turns
them into a five-minute model of an inference deployment --- step time, the
batch size at which decode stops being bandwidth-bound, tokens per second per
user and in aggregate, time to first token, memory headroom and dollars per
million tokens at a latency target. Each prediction is checked against a
measurement, and the chapter catalogues where napkin math goes wrong: kernel
efficiency, launch and scheduling overhead, communication, and real workload
mixes.
\end{ChapterIntent}

\OpeningSection{Predicting a benchmark before running it}

\NapkinSection{Napkin math: the step-time model}

\section{Batch size: the throughput--latency curve}

\section{The knee: where decode turns compute-bound}

\section{Adding the KV cache}

\section{Prefill in the same model}

\section{From tokens per second to dollars per million tokens}

\section{Goodput: throughput that meets its objective}

\section{Worked examples on real hardware}

\section{Where napkin math breaks}

\LedgerSection

\LabSection{a cost calculator you can trust}

\HappenedSection{predictions against a published system}

\FrontierSection{}

\ExercisesSection
